\SourceLeaf{417}%
für den betreffenden singulären Fall gibt es tatsächlich immer vier Lösungen. Man legt entweder durch \((4)\) oder durch \((5)\) eine der beiden Geraden, deren Winkel mit der Abszissenachse die Größe \(\frac A{2a}=\frac B{2b}\) als Kosinus hat. Ob diese Lösungen reell sind, hängt davon ab, dass diese Größe nicht größer als \(1\) ist. Bei genau diesem Wert reduzieren sich die vier Lösungen auf zwei.
\par\bigskip\begin{center}\rule{.16\textwidth}{.3pt}\end{center}
